% =====================================================================
\section{Passive-regime witness: b-anomaly}\label{sec:passive_witness}
% =====================================================================

This section threads the first of the two domain-disjoint empirical landings of
$\Cph$. The full preprint is~\citep{SmartBAnomaly2026}; we
summarise here only what the operator-witness narrative requires
and inherit the preprint's caveats verbatim.

\subsection{What b-anomaly tests}\label{ssec:banomaly_setup}

The Wilson-coefficient Hamiltonian for $b\to s\mu^{+}\mu^{-}$
contains a $C_{9}^{(\prime)}$ contribution that, in the Standard
Model, is approximately $q^{2}$-independent in the relevant
kinematic range. The b-anomaly preprint tests a fixed
$q^{2}$-dependent effective shift on top of the SM backend, of the
form
\begin{equation}\label{eq:banomaly_kernel}
\Delta C_{9}^{\mathrm{eff}}(q^{2})
\;=\;
-\,A \cdot \kappa_{V_{600}}(q^{2}),
\end{equation}
where $\kappa_{V_{600}}(q^{2}) > 0$ on the LHCb $q^{2}$ window is
the projection of $\Cph$ on $\Rsixhundred$ to the flavour-physics
$q^{2}$ axis (the b-anomaly preprint's §3 projection construction,
which we do not relitigate here; this is a projection of the
operator, not a derivation of $\Ph^{-2}$), and $A$ is a single
fitted dimensionless amplitude per dataset. The explicit minus
sign in Eq.~\eqref{eq:banomaly_kernel} reflects the b-anomaly
preprint's sign convention: a positive fitted $A>0$ produces a
negative $\Delta C_{9}^{\mathrm{eff}}<0$, the established
direction of the angular anomaly. The kernel shape
$\kappa_{V_{600}}$ is held fixed across all five datasets.
This is a \emph{structural} test: same fixed $\Cph$ on the same
$\Rsixhundred$, no shape retuning between datasets.

\subsection{The five-dataset structural fit}

The b-anomaly preprint reports the following per-dataset table
(verbatim from~\citep{SmartBAnomaly2026}, also at
\texttt{BANOMALY-001/vfd-b-anomaly/README.md}):

\begin{table}[ht]
\centering
\small
\caption{b-anomaly five-dataset structural fit. Verbatim
from~\citep{SmartBAnomaly2026}; one fitted amplitude $A$ per
dataset, kernel shape held fixed.}
\label{tab:banomaly}
\begin{tabular}{l l r r r r}
\toprule
Dataset & Decay & $n$ & $\Delta\mathrm{AIC}_{\mathrm{NL}}$ &
   Best-fit $A$ & $\Delta C_{9}^{\mathrm{eff}}$ \\
\midrule
LHCb 2015 & $B^{0}\!\to\!K^{*0}$ & $32$ & $-0.24$ & $+1.24$ & $-0.96$ \\
LHCb 2021 & $B^{+}\!\to\!K^{*+}$ & $32$ & $+0.17$ & $+2.06$ & $-1.59$ \\
CMS 2025 (no $P_{4}'$) & $B^{0}\!\to\!K^{*0}$ & $18$ & $+0.47$ & $+1.05$ & $-0.81$ \\
LHCb 2025 & $B^{0}\!\to\!K^{*0}$ & $32$ & $+1.09$ & $+1.14$ & $-0.86$ \\
LHCb 2015 & $B_{s}\!\to\!\phi$ ($S$-basis) & $24$ & $-0.24$ & $+4.98$ & $-3.85$ \\
\bottomrule
\end{tabular}
\end{table}

\subsection{What the structural fit reports}

\begin{itemize}\itemsep=2pt
\item \textbf{Cross-dataset consistency (5/5).} The same fixed
  kernel shape can be fit to all five datasets with one amplitude
  $A$ per dataset and no shape retuning. The kernel never moves
  between datasets.
\item \textbf{Sign uniformity (5/5).} $A > 0$ in $5/5$ fits;
  $\Delta C_{9}^{\mathrm{eff}} < 0$ in $5/5$ fits. The kernel
  reproduces the established direction of the
  anomaly~\citep{LHCbAngular2020} across all five independent
  measurements.
\item \textbf{Cross-channel ratio.} The b-anomaly README reports
  the basis amplification (the predicted Krüger--Matias $P$-basis
  vs $S$-basis factor $\sim 2.2$~\citep{KrugerMatias2005}) as
  partly explanatory for the $B\to K^{*}$ vs $B_{s}\!\to\!\phi$
  amplitude gap. Working through the explicit arithmetic with
  $B\to K^{*}$ best-fit $A_{P} \approx 1.14$: the basis-corrected
  prediction for $B_{s}\!\to\!\phi$ is
  $A_{P} \cdot 2.2 \approx 2.5$; the observed $B_{s}\!\to\!\phi$
  amplitude is $A \approx 4.98$, leaving an unresolved amplitude
  excess of about a factor two above the basis-corrected
  prediction. The b-anomaly preprint reports this residual as an
  open item, not a discharge.
\item \textbf{Geometry-first variant test.} Of three discrete
  Laplacian variants on $\Rsixhundred$ (unweighted,
  $\Ph$-geometric weighted, $\Ph$-arithmetic weighted), the
  unweighted choice ranks highest on both a pure-geometry
  criterion (correlation $0.997$ with the continuum kernel under
  the b-anomaly preprint's $q^{2}$/shell-projection geometry
  metric, §3.4 of~\citep{SmartBAnomaly2026} — \emph{not} this
  paper's per-vertex test, whose unweighted score is $0.976$;
  the two metrics differ in what is correlated) and the LHCb~2025
  data $\chi^{2}$ ($\chi^{2}=13.555$). The two criteria agree on
  the variant ranking — a two-criterion agreement on the same
  fixed operator.
\end{itemize}

\subsection{What the structural fit does \emph{not} establish}

The b-anomaly preprint is explicit about the following caveats,
which we inherit verbatim:

\begin{itemize}\itemsep=2pt
\item \textbf{AIC tie on current data.} On Akaike model comparison,
  $\Cph$-derived $\kappa_{V_{600}}$ and a constant Wilson-coefficient
  shift ($\mathrm{FREE\_C9}$, also $k=1$) are statistically
  indistinguishable: stacked AIC weights
  $w_{\mathrm{VFD}} = 0.348$ vs $w_{\mathrm{FREE\_C9}} = 0.652$.
  Current data cannot resolve the model comparison. AIC is blind
  to the universality / shape-prediction claim itself, but it is
  decisive about whether the shape is forced by data: it is not.
\item \textbf{Free-width Gaussian alternative.} A free-width
  Gaussian charm-loop proxy fits the same five datasets comparably
  in $\chi^{2}$ at the cost of one extra shape parameter; $\Cph$
  is not the unique $q^{2}$ shape consistent with the anomaly.
\item \textbf{Mode-B drift (linearised-to-non-linear refit).} An
  earlier analysis the b-anomaly project labels Mode-B
  (linearised) gave a stronger preference
  ($\Delta\mathrm{AIC} = -1.67$ on LHCb 2025) that did not survive
  the subsequent non-linear refit; the $+2.77$-AIC-unit drift
  between Mode-B (linearised) and the non-linear refit is the
  largest single methodological uncertainty in the b-anomaly
  project.
\item \textbf{Look-elsewhere on the variant test.} The b-anomaly
  preprint's limitations section~\citep{SmartBAnomaly2026}
  acknowledges that the LHCb~2025 data was looked at first, and
  only later was the agreement of the data-$\chi^{2}$ ranking with
  the pure-geometry ranking verified. (The b-anomaly README
  emphasises that the geometry-only criterion is independent of
  the LHCb data; the historical non-blindness is in the project's
  ordering, not in the criterion.) The two-criterion agreement is
  criterion-independent but not historically blind.
\end{itemize}

\subsection{Reading at the operator level}

The b-anomaly result is the \emph{passive-regime} empirical
witness for $\Cph$ on $\Rsixhundred$: a single linear response
$\psi = \Cph^{-1} f$, projected to the $q^{2}$ axis through a
fixed discrete-to-momentum projection, gives a sign-uniform
structural fit consistent with the $b\to s\mu^{+}\mu^{-}$ angular
anomaly across five independent measurements without shape
retuning. This does
not establish the kernel as theorem-grade physics on the flavour
side (the AIC tie, the free-width Gaussian alternative, and the
Mode-B linearised-to-non-linear refit drift prevent that). It
does support, at the operator-witness level, the inherited reading
that the same fixed $\Cph$ on the same fixed $\Rsixhundred$ is
consistent with one of two domain-disjoint empirical landings without
parameter retuning. The
second landing is in \S\ref{sec:active_witness}.
